% !PRINTED 215
\weberpage{215}
%<<CONT

\weq{16}{(\gamma\xi+\delta)^{n-2}\, T\!\left(\frac{\alpha\xi+\beta}{\gamma\xi+\delta}\right) = r^{n-2}\Theta(\xi),}

worin der Zusammenhang zwischen den Co\"efficienten $t$ und $\tau$
durch die symbolischen Gleichungen (13) ausgedr\"uckt, also derselbe
 ist, wie in den Functionen $Y$ und $H$.

\gesperrt{Es ist also hiernach $r^{n-2}\Theta(\xi)$ die Umformung der
Function $(\gamma\xi+\delta)^{n-2}T(z)$ durch dieselbe lineare Substitution,
durch die $(\gamma\xi+\delta)^n f(z)$ in $\varphi(\xi)$ \"ubergeht.}

Bezeichnen wir die Co\"efficienten dieser umgeformten Function,
also die Gr\"ossen $r^{n-2}\tau_k$ mit $t'$, so ergiebt die Gleichung (12),
die in Bezug auf $\tau$ linear ist,

\weq{17}{H(t', \xi) = r^{n-1}Y(t, x).}

Aus $Y(t, x)$ gehen nun $n$ Functionen hervor, wenn man
f\"ur $x$ die $n$ Wurzeln von $f(x)$ setzt. Jede homogene, rationale,
symmetrische Function dieser $n$ Functionen ist rational durch
die Co\"efficienten von $f(x)$ ausdr\"uckbar, und wenn wir also eine
solche Function mit $K(t, a)$ bezeichnen, so ergiebt die Gleichung
(17), wenn wir die Co\"efficienten von $\varphi(\xi)$ mit $a'$ bezeichnen,

\weq{18}{K(t', a') = r^{\nu(n-1)}K(t, a),}

wenn $\nu$ den Grad der symmetrischen Function bedeutet. Damit
ist der sch\"one Satz von \gesperrt{Hermite} bewiesen:

\gesperrt{Die Co\"efficienten in der durch die Tschirnhausen-
Transformation}

\[
y = Y(t, x)
\]

\gesperrt{umgeformten Gleichung $f(x) = 0$ und alle symmetrischen
 Functionen der $n$ Werthe $y$ sind simultane
Invarianten der beiden Functionen}

\[
f(z), \quad T(z).
\]

\begin{center}
\S.\,70.\\[2pt]
\gesperrt{Ausf\"uhrungen \"uber den Hermite'schen Satz.}
\end{center}

Durch die Betrachtungen des letzten Paragraphen hat sich
ergeben, dass, wenn wir

\weq{1}{y = t_{n-2}F_0(x) + t_{n-3}F_1(x) \cdots + t_0 F_{n-2}(x)}

setzen, die symmetrischen Functionen der $n$ Werthe $y_1, y_2 \cdots y_n$,
die man aus $y$ erh\"alt, wenn man f\"ur $x$ die $n$ Wurzeln von
%CONT>>
